\section{Results}\label{sec:results}
\begin{table}[t]
\centering\scriptsize
\setlength{\tabcolsep}{3.2pt}
\caption{Fresh-seed results (seeds 100--119, mean$\pm$s.d.; frozen protocol). GT dist.: ground-truth algebra
distance after Procrustes alignment (chance $0.80$ for $T^2$, $0.50$ for SO(3)). CC: correct-closed runs. SO(2), where
all methods coincide, is in Appendix Table~\ref{tab:mainfull}.}
\label{tab:main}
\begin{tabular}{llccccc}
\toprule
Group & Method & 1-step MSE & 10-step MSE & Closure residual & GT dist. & CC\\
\midrule
\tabinput{../../results/review/tables/main_t2so3.tex}
\bottomrule
\end{tabular}
\end{table}

\begin{table}[t]
\centering\scriptsize
\setlength{\tabcolsep}{2.6pt}
\caption{Pre-registered primary hypotheses, BracketNet vs.\ \textsc{+Comp.} (one-sided). Units: seeds ($n{=}20$)
for H1--H2, disjoint seed pairs ($n{=}10$) for H3. Reported: two-sided 95\% $t$-interval of the difference, paired
effect size $d_z$, and the number of units where BracketNet is lower. $p$: exact sign-flip; Holm over six tests.}
\label{tab:primary}
\begin{tabular}{llcccccccc}
\toprule
Hypothesis & Group & $n$ & +Comp. & BracketNet & 95\% CI ($t$) & $d_z$ & lower & $p$ & Holm $p$\\
\midrule
\tabinput{../../results/final/tables/primary.tex}
\bottomrule
\end{tabular}
\end{table}

\subsection{Closure makes SO(3) algebras closed, but often wrong}
\paragraph{Transition fitting closes $T^2$ but not SO(3).}
Without closure regularization, the true $T^2$ algebra is recovered in \NTtwoCompCatCorrectclosed/20 (+Comp.) and
\NTtwoLocalCatCorrectclosed/20 (Local) runs, while the SO(3) span stays \emph{non-closed} in
\NSOthreeCompCatNonclosed/20 and \NSOthreeLocalCatNonclosed/20 runs. One- and ten-step transport errors are low in
both cases (Table~\ref{tab:main}). BracketNet at the originally planned $\lambda{=}0.1$ behaves like the baselines.

\paragraph{Closure regularization: large effects for SO(3), concentrated effects for $T^2$.}
For SO(3), closure falls from \NHoneSOthreeMeanb{} to \NHoneSOthreeMeana{} (lower on \NHoneSOthreeWins{} seeds), and
the ground-truth distance falls from \NHtwoSOthreeMeanb{} to \NHtwoSOthreeMeana{} (\NHtwoSOthreeWins{}). Both are
significant after Holm correction (Table~\ref{tab:primary}).

For $T^2$ the closure test is also significant, but the effect is concentrated: the median paired difference is
\RInvHoneTtwoMedian{}, and the differences are heavily skewed (skewness \RInvHoneTtwoSkew). The $t$-interval
includes zero, while the interval obtained by inverting the primary sign-flip test,
$[\RInvHoneTtwoLo,\RInvHoneTtwoHi]$, does not. The difference comes from the few seeds in which \textsc{+Comp.}\ fails
to close. The $T^2$ ground-truth distance does not change significantly.

Ten-step transport shows no detectable difference in either group; the 95\% intervals are
$[\NTtwoBNVsMseonezerostepCilo,\NTtwoBNVsMseonezerostepCihi]$ and
$[\NSOthreeBNVsMseonezerostepCilo,\NSOthreeBNVsMseonezerostepCihi]$.

\paragraph{Three SO(3) outcomes.}
Every run falls into one of the classes of Prop.~\ref{prop:classify} or is non-closed (Fig.~\ref{fig:mech}a).
\begin{itemize}
\item BracketNet: \NSOthreeBNCatCorrectclosed{} correct (diagonal), \NSOthreeBNChiral{} chiral, \NSOthreeBNCatNonclosed{} non-closed.
\item \textsc{+Comp.}: \NSOthreeCompCatCorrectclosed{} correct, \NSOthreeCompCatIncorrectclosed{} chiral, \NSOthreeCompCatNonclosed{} non-closed.
\end{itemize}
Closure converts non-closed runs into closed ones, and about half of them land in the wrong class. The rise in
correct-closed runs is not significant (exact McNemar $p=\NSOthreeMcnemarp$).

\subsection{Mechanism: the endpoint objective favours chiral solutions}\label{sec:mechanism}
We tested Prop.~\ref{prop:endpoint} with predictions registered before any of the following runs.

\paragraph{Composition error by class (existing runs).}
As predicted, chiral runs compose far better than correct ones: median composition error \REoneIncorrectclosedchiralComp{}
(chiral) vs.\ \REoneCorrectclosedComp{} (correct). Their one-step errors are the same (\REoneIncorrectclosedchiralMseone{}
vs.\ \REoneCorrectclosedMseone). Their action disagreement with the truth is \REoneIncorrectclosedchiralAction{} vs.\
\REoneCorrectclosedAction. The correct runs' value is the floor imposed by Prop.~\ref{prop:endpoint}(c); chiral
runs are no better than the identity.

\paragraph{Closure geometry alone favours the correct class.}
Pure closure flow from random subspaces reaches the diagonal class in \RFlowDiag/\RFlowN{} starts
(Sec.~\ref{sec:theory}), so the high chiral rate in training is not a property of the closure penalty alone.

\paragraph{Trace of the generator span.}
We re-ran the 20 BracketNet SO(3) runs with generator snapshots; \REtwoIdentical{} reproduced the stored runs
bit-for-bit. We measured each span's chirality $\chi=e_L-e_R$, the share of its energy in $\su(2)_L$ minus that
in $\su(2)_R$.
\begin{itemize}
\item \emph{Pre-registered.} At the start of the closure ramp, $|\chi|$ predicts a chiral outcome among the
\REtwoNclosed{} closed runs with AUC \REtwoAucramp{} (at initialization: \REtwoAucinit). Medians are
\REtwoChirampIncorrectclosedchiral{} (ending chiral) vs.\ \REtwoChirampCorrectclosed{} (ending correct).
\item \emph{Post hoc.} The closure residual at the ramp start separates outcomes more sharply (AUC
\REtwoAucclosure). \REtwoCorrectnear{} of the \REtwoNcorrect{} runs that end correct had already nearly closed onto
the diagonal during the transition-fitting phase. \REtwoChiralfar{} of the \REtwoNchiral{} chiral runs were still
far from closed when the penalty switched on (Fig.~\ref{fig:mech}b).
\end{itemize}
Together with the closure flow, this suggests the following picture. Transition fitting sometimes finds the
diagonal algebra by itself, and closure then keeps it. When the span is still far from closed, closing it under the
transport and composition losses, which favour chiral solutions at the ground-truth embedding
(Prop.~\ref{prop:endpoint}), usually yields a
chiral algebra.

\paragraph{Intervention: supplying the true action coefficients (pre-registered).}
Prop.~\ref{prop:endpoint} predicts that chiral solutions lose their advantage if actions need not be inferred from
endpoints. On 20 new seeds (400--419) we trained BracketNet with the true coefficients in the transport terms. This
is the action-supervised setting of homomorphism autoencoders \citep{keurti2023homomorphism}; $q$ is fit only for
evaluation. Results, relative to standard BracketNet on the same seeds:
\begin{itemize}
\item chiral outcomes fell from \REthreeBracketnetChiral/20 to \REthreeOraclebnChiral/20 (exact McNemar $p=\REthreeChiralp$);
\item correct-closed runs rose from \REthreeBracketnetCorrect/20 to \REthreeOraclebnCorrect/20 ($p=\REthreeCorrectp$);
\item the ground-truth distance changed by $\REthreeGtdiff$ ($[\REthreeGtlo,\REthreeGthi]$);
\item without any closure term, the oracle model also reaches \REthreeOraclelocalCorrect/20 correct-closed runs.
\end{itemize}
Once actions are given, closure regularization is unnecessary: the difficulty lies in inferring actions from
endpoints. The oracle models' own $q$, which must infer actions from endpoints at test time, still has action
disagreement \REthreeOraclebnAction, the floor of Prop.~\ref{prop:endpoint}(c). The intervention adds
information as well as removing the ambiguity, so it supports the mechanism but does not isolate it.

\begin{figure}[t]\centering
\includegraphics[width=\linewidth]{fig_mechanism.pdf}
\caption{(a) Final SO(3) span chirality $\chi$ vs.\ closure residual for every fresh-seed run: closed runs sit at
$\chi\in\{-1,0,1\}$, the classes of Prop.~\ref{prop:classify}. (b) Closure of the BracketNet span when the closure
penalty switches on, by final outcome (post hoc). (c) Composition error by outcome class, all SO(3) runs.
(d) Outcomes with endpoint-inferred vs.\ true (oracle) action coefficients, seeds 400--419.}\label{fig:mech}
\end{figure}

\subsection{A rendered-image benchmark}\label{sec:render}
To test a harder observation process, we rendered the SO(3) state as $20{\times}20$ images: the 3-D position of a
Gaussian blob sets its image position and size, and the invariant coordinate sets its brightness ($p=400$). The
frozen protocol was applied unchanged on 20 new seeds (300--319), with predictions registered in advance.
\begin{itemize}
\item \emph{No baseline run closes} (Local \REfourLocalNonclosed/20, \textsc{+Comp.} \REfourCompNonclosed/20 non-closed).
BracketNet lowers closure on \REfourClosurebnvscompWins{} seeds (\REfourBracketnetClosure{} vs.\ \REfourCompClosure,
one-sided $p=\REfourClosurebnvscompP$).
\item \emph{Every closed BracketNet run is chiral} (\REfourBracketnetChiral/20; none diagonal). This is read off the
chirality invariant, which needs no latent alignment.
\item \emph{Ground-truth recovery cannot be assessed by our alignment metric here.} The learned latents are far from
an isometry of the true state (CKA with it \REfourBracketnetCkatruth, Procrustes residual \REfourBracketnetProcrustestruth).
\item \emph{There is a small transport cost}: 10-step error rises by $\REfourMseonezerobnvscompMeandiff$
($[\REfourMseonezerobnvscompCilo,\REfourMseonezerobnvscompCihi]$).
\end{itemize}
Under rendering, closure regularization therefore produced closed algebras only of the wrong class. The benchmark is
still synthetic and does not test natural images.
